\chapter{Knee Replacement}

\section*{What Is This Surgery?}
Knee replacement, medically termed total knee arthroplasty (TKA), is a major reconstructive surgical procedure involving the resurfacing of the severely damaged weight-bearing surfaces of the knee joint. This intervention primarily addresses end-stage degenerative conditions, most commonly osteoarthritis, but also rheumatoid arthritis and post-traumatic arthritis, that have led to debilitating pain, significant loss of function, and progressive deformity. During the procedure, the worn-out articular cartilage and a small amount of underlying subchondral bone from the distal femur, proximal tibia, and often the patella are precisely resected and replaced with engineered prosthetic components. These components, typically made of metal alloys for the femoral and tibial ends and a high-density polyethylene spacer, are designed to restore smooth articulation and mechanical alignment. The overarching goal of TKA is to alleviate chronic knee pain, correct existing deformities, and restore functional mobility, thereby profoundly enhancing the patient's quality of life and independence.

\section*{Alternative Names}
\begin{itemize}
  \item Total Knee Arthroplasty (TKA)
  \item Knee Reconstruction Surgery
  \item Knee Joint Replacement
  \item Arthroplasty of the Knee
  \item Total Condylar Arthroplasty
  \item Unicompartmental Knee Arthroplasty (Partial Knee Replacement) - *Note: While TKA is the focus, UKA is a related alternative.*
\end{itemize}

\section*{Relevant Surgical Anatomy}
The knee joint is a complex synovial hinge joint, primarily formed by the articulation of the distal femur and the proximal tibia, with the patella (kneecap) articulating anteriorly with the femoral trochlea. The articular surfaces of the femoral condyles, tibial plateau, and posterior patella are normally covered by smooth hyaline cartilage, which facilitates low-friction movement. The medial and lateral menisci, C-shaped fibrocartilaginous structures, sit between the femoral and tibial condyles, acting as shock absorbers and enhancing joint congruity. Stability is provided by a robust ligamentous complex, including the medial collateral ligament (MCL) and lateral collateral ligament (LCL) which prevent varus and valgus stress, respectively, and the anterior cruciate ligament (ACL) and posterior cruciate ligament (PCL) which control anterior and posterior tibial translation.

The quadriceps femoris muscle, comprising four heads, converges into the quadriceps tendon, which inserts into the patella and continues as the patellar tendon to the tibial tuberosity, forming the extensor mechanism crucial for knee extension. Neurovascular structures of critical importance during TKA include the popliteal artery and vein, which lie posteriorly within the popliteal fossa, and the common peroneal nerve, which courses laterally around the fibular head. Damage to these structures can lead to devastating complications. Lymphatic drainage of the knee region primarily follows the major vascular bundles to the popliteal and then inguinal lymph nodes. Surgical landmarks include the medial and lateral epicondyles of the femur, the tibial tuberosity, and the patella itself, which guide incision placement and bone resection.

\section{Surgical Principle \& Rationale}
The fundamental surgical principle of total knee arthroplasty is to eliminate pain and restore functional mobility by precisely replacing diseased or damaged articular surfaces with artificial implants. In conditions such as severe osteoarthritis, the protective articular cartilage erodes, leading to painful bone-on-bone friction, chronic inflammation, and progressive joint destruction, often accompanied by significant deformity (e.g., varus or valgus malalignment). The rationale for TKA is that by removing these damaged surfaces and replacing them with smooth, biocompatible prosthetic components, the source of pain is eliminated, and the mechanical axis of the limb is restored, allowing for pain-free movement.

The operative concept involves meticulous bone resection to create precise beds for the implants, followed by careful soft tissue balancing to ensure joint stability throughout the full range of motion. The technical goals include restoring the anatomical and mechanical axes of the lower limb, achieving a balanced flexion and extension gap, ensuring adequate ligamentous tension, and securing durable fixation of the implants to the host bone, typically with bone cement. The femoral component resurfaces the distal femur, the tibial component covers the proximal tibia, and a polyethylene insert acts as the new bearing surface. Patellar resurfacing, when performed, replaces the posterior surface of the kneecap with a polyethylene button. This comprehensive reconstruction aims to provide a long-lasting, stable, and mobile knee joint, significantly improving the patient's quality of life.

\section{History \& Surgical Evolution}
The journey of knee arthroplasty began in the late 19th and early 20th centuries with rudimentary attempts at interpositional arthroplasty, using materials like fascia, fat, or even pig bladder to create a new joint surface. However, these early efforts yielded inconsistent results and were prone to failure. The true dawn of modern total knee replacement surgery emerged in the 1960s and 1970s, driven by advancements in biomaterials and surgical understanding.

A significant milestone was the development of the Geomedic knee by Frank Gunston in 1968, which introduced the concept of resurfacing the femoral condyles and tibial plateau separately, moving away from the highly constrained hinge designs that were prone to loosening. The most pivotal breakthrough came with the introduction of the Total Condylar Knee by John Insall, Peter Walker, and others at the Hospital for Special Surgery in 1974. This revolutionary design featured a metal femoral component, a metal-backed tibial component with a polyethylene insert, and often a polyethylene patellar component, establishing the fundamental design principles that largely persist today. Subsequent decades saw continuous evolution, including improvements in implant materials (e.g., highly cross-linked polyethylene for reduced wear), fixation methods (cemented versus cementless options), and design variations such as posterior-stabilized (PS) and cruciate-retaining (CR) prostheses to better mimic natural knee kinematics. More recently, computer-assisted navigation, robotic-assisted surgery, and patient-specific instrumentation have been introduced to enhance precision in bone cuts and implant alignment, further improving outcomes and longevity.

\section{Clinical Indications}
Total knee arthroplasty is indicated for patients experiencing severe, debilitating knee pain and functional impairment due to end-stage knee joint disease that has failed to respond to comprehensive non-surgical management.

\begin{itemize}
  \item Severe, chronic knee pain that significantly limits daily activities, sleep, and quality of life, unresponsive to conservative treatments such as physical therapy, anti-inflammatory medications, and injections.
  \item Radiographic evidence of advanced degenerative joint disease, including severe osteoarthritis (primary or secondary), rheumatoid arthritis, or post-traumatic arthritis, characterized by bone-on-bone articulation, significant joint space narrowing, osteophytes, and subchondral sclerosis.
  \item Significant functional impairment, such as difficulty walking, climbing stairs, standing for prolonged periods, or performing activities of daily living, often requiring assistive devices.
  \item Progressive knee deformity, including varus (bow-legged) or valgus (knock-kneed) malalignment, which contributes to pain, instability, and abnormal gait mechanics.
  \item Persistent knee stiffness or a severe loss of range of motion that interferes with normal function and hygiene.
  \item Failure of previous knee surgeries, such as arthroscopy, osteotomy, or unicompartmental knee arthroplasty, to provide lasting pain relief or functional improvement.
  \item Inflammatory arthropathies (e.g., psoriatic arthritis, ankylosing spondylitis) leading to severe, destructive joint changes.
  \item Avascular necrosis (osteonecrosis) of the femoral condyles causing severe pain, joint collapse, and functional limitation.
\end{itemize}

\section{Clinical Positioning}
Total knee arthroplasty is positioned as a definitive, reconstructive surgical treatment for end-stage knee joint disease. It is generally not a first-line intervention but rather a secondary or tertiary option pursued only after a comprehensive trial of conservative, non-surgical treatments has failed to provide adequate pain relief or functional improvement. These conservative measures typically include activity modification, weight management, physical therapy, oral analgesics and anti-inflammatory medications, corticosteroid injections, and viscosupplementation.

TKA becomes the primary definitive treatment when the severity of joint degeneration, intractable pain, and profound functional limitation reaches a point where only surgical reconstruction can offer substantial and lasting relief. In rare, acute scenarios, such as severe trauma with irreparable joint damage or rapidly progressing destructive arthritis, TKA might be considered earlier in the treatment algorithm. However, for the vast majority of chronic degenerative conditions, it represents the final, most effective step in a progressive treatment pathway, aimed at restoring function and significantly improving quality of life when all other options have been exhausted.

\section{Surgical Technique \& Procedure}
Total knee arthroplasty is performed in a meticulously sterile operating room environment, typically lasting between 1 to 2 hours, depending on the complexity of the case.

\begin{itemize}
  \item \textbf{Anesthesia:} The procedure is most commonly performed under general anesthesia, spinal anesthesia, or a combination of spinal and regional nerve blocks (e.g., adductor canal block, femoral nerve block) to provide effective perioperative pain control and minimize systemic opioid requirements.
  \item \textbf{Patient Positioning \& Preparation:} The patient is positioned supine on the operating table. The affected leg is meticulously prepared and draped in a sterile fashion. A pneumatic tourniquet is often applied to the upper thigh and inflated to achieve a bloodless field, which enhances visibility and precision during bone cuts and cement application.
  \item \textbf{Incision:} A longitudinal skin incision, typically 15-25 cm in length, is made over the anterior aspect of the knee. The most common approach is the medial parapatellar arthrotomy, which involves incising the skin, subcutaneous tissue, and then the joint capsule medial to the patella, allowing for eversion of the patella and full exposure of the femoral and tibial condyles.
  \item \textbf{Bone Resection:} Specialized cutting guides, often patient-specific or navigated, are used to precisely resect the damaged articular surfaces of the distal femur, proximal tibia, and, if indicated, the posterior surface of the patella. These cuts are critical for restoring proper limb alignment, joint line, and creating precise bone beds for the prosthetic components.
  \item \textbf{Soft Tissue Balancing \& Trial Implants:} After bone preparation, trial prosthetic components are inserted. The surgeon then meticulously assesses joint stability, ligamentous balance in both flexion and extension, and the knee's range of motion. Adjustments, such as controlled soft tissue releases (e.g., posterior capsular release, lateral retinacular release), are performed as needed to achieve optimal balance and kinematics.
  \item \textbf{Definitive Implantation:} Once optimal fit, alignment, and balance are achieved with the trial components, the definitive metal femoral component, metal tibial tray, and polyethylene insert are cemented into place using bone cement (polymethylmethacrylate). If patellar resurfacing is performed, a polyethylene patellar button is also cemented to the prepared posterior surface of the kneecap.
  \item \textbf{Closure:} The joint is thoroughly irrigated, and the tourniquet is deflated. Hemostasis is achieved. The surgical layers are meticulously closed, including the joint capsule, muscle, subcutaneous tissue, and skin, often with absorbable sutures and skin staples. A drain may be placed temporarily to collect any excess fluid or blood. A sterile dressing is applied.
\end{itemize}

\section*{Preoperative Workup \& Preparation}
Thorough preoperative workup and preparation are paramount for optimizing patient outcomes, minimizing surgical risks, and ensuring a smooth recovery from total knee arthroplasty.

\begin{itemize}
  \item \textbf{Comprehensive Medical Evaluation:} A detailed medical evaluation by the primary care physician and/or an anesthesiologist is mandatory to assess the patient's overall health status, identify any significant medical comorbidities (e.g., cardiovascular disease, diabetes, pulmonary conditions, renal insufficiency), and ensure the patient is medically optimized and fit for surgery. This often includes:
    \begin{itemize}
      \item Blood tests: Complete blood count (CBC), basic metabolic panel (BMP), liver function tests (LFTs), coagulation profile (PT/INR, PTT), and type and screen.
      \item Urinalysis: To rule out urinary tract infections.
      \item Electrocardiogram (ECG): To assess cardiac rhythm and function.
      \item Chest X-ray: To evaluate lung status.
      \item Cardiac clearance: For patients with significant cardiac history, further evaluation by a cardiologist may be required.
    \end{itemize}
  \item \textbf{Medication Review \& Adjustment:} All medications, particularly anticoagulants (e.g., warfarin, DOACs) and antiplatelet agents (e.g., aspirin, clopidogrel), must be meticulously reviewed. Patients are typically instructed to stop these medications several days to a week prior to surgery, under the strict guidance of their prescribing physician, to minimize bleeding risk.
  \item \textbf{Dental Clearance:} A recent dental examination and clearance are often required to rule out any active dental infections or significant periodontal disease, which could potentially serve as a source of bacteremia and lead to periprosthetic joint infection.
  \item \textbf{Physical Therapy/Pre-habilitation:} Some patients benefit from "pre-habilitation" with physical therapy prior to surgery. This focuses on strengthening the quadriceps and hamstring muscles, improving knee range of motion, and educating the patient on postoperative exercises and mobility techniques.
  \item \textbf{NPO Status:} Patients are strictly instructed to be NPO (nil per os - nothing by mouth) for a specified period (typically 6-8 hours) before surgery to prevent aspiration during anesthesia induction.
  \item \textbf{Prophylactic Antibiotics:} Intravenous prophylactic antibiotics (e.g., cefazolin) are administered within 60 minutes prior to the surgical incision to significantly reduce the risk of surgical site infection.
  \item \textbf{Informed Consent:} The patient must fully understand the surgical procedure, its potential benefits, risks, alternatives, and expected recovery course, and provide comprehensive informed consent.
  \item \textbf{Home Preparation:} Patients are advised to prepare their home environment for recovery, including removing tripping hazards, arranging for assistance with daily tasks, and setting up a comfortable recovery area with necessary aids (e.g., elevated toilet seat, shower chair).
\end{itemize}

\section*{Contraindications \& Precautions}
Careful patient selection is crucial for successful total knee arthroplasty. Certain conditions may preclude surgery or necessitate extreme caution.

\textbf{Absolute Contraindications:}
\begin{itemize}
  \item Active infection in the knee joint (septic arthritis) or elsewhere in the body (e.g., urinary tract infection, dental abscess, skin infection), as this dramatically increases the risk of devastating periprosthetic joint infection.
  \item Severe peripheral vascular disease with compromised circulation to the affected limb, which can severely impair wound healing, increase infection risk, and lead to limb-threatening complications.
  \item Dysfunction of the extensor mechanism (e.g., quadriceps rupture, severe patellar tendon pathology) that would prevent effective postoperative rehabilitation and ambulation.
  \item Neuropathic arthropathy (Charcot joint) with severe bone destruction, joint instability, and impaired pain sensation, which can lead to rapid implant failure.
  \item Skeletally immature patients with open growth plates, as TKA would disrupt normal bone growth.
  \item Uncontrolled medical comorbidities that pose an unacceptably high anesthetic or surgical risk, such as recent myocardial infarction (within 3-6 months), unstable angina, severe uncontrolled diabetes (HbA1c >8.0-8.5\%), or severe pulmonary hypertension.
  \item Active osteomyelitis in the affected limb.
\end{itemize}

\textbf{Relative Contraindications \& Precautions:}
\begin{itemize}
  \item Morbid obesity (BMI >40 kg/m\textasciicircum{}2), which significantly increases surgical complexity, implant stress, and the risks of complications including infection, deep vein thrombosis (DVT), pulmonary embolism (PE), wound healing issues, and early implant failure. Weight loss is often recommended prior to surgery.
  \item Severe osteoporosis, which may compromise implant fixation and increase the risk of periprosthetic fracture. Bone density optimization may be required.
  \item Poor skin condition around the knee (e.g., active psoriasis, chronic venous stasis ulcers, severe dermatitis), increasing the risk of surgical site infection and wound complications.
  \item Significant psychiatric illness, cognitive impairment, or substance abuse that may hinder cooperation with the rigorous postoperative rehabilitation protocol.
  \item Unrealistic patient expectations regarding pain relief, functional recovery, or return to high-impact activities.
  \item Severe flexion contracture or ankylosis (fusion) of the knee, which may make achieving a functional range of motion challenging and increase the risk of stiffness.
  \item History of prior knee infection, even if successfully treated, as it carries a higher lifetime risk of recurrence and periprosthetic joint infection.
  \item Chronic pain syndromes or fibromyalgia, which may complicate postoperative pain management and patient satisfaction.
\end{itemize}

\section*{Complications \& Risks}
While total knee replacement is a highly successful and generally safe procedure, like any major surgery, it carries potential risks and complications that patients must be aware of.

\textbf{General Surgical Complications:}
\begin{itemize}
  \item \textbf{Anesthesia-related risks:} These include adverse reactions to anesthetic agents, respiratory depression, cardiac events (e.g., myocardial infarction, arrhythmia), stroke, and nerve damage from regional blocks (e.g., transient or permanent numbness/weakness).
  \item \textbf{Bleeding:} Excessive blood loss during or after surgery may necessitate blood transfusion, which carries its own risks (e.g., transfusion reactions, infection). Hematoma formation can also occur, potentially requiring drainage.
  \item \textbf{Infection:} Surgical site infection can range from superficial wound infections (cellulitis) to devastating deep periprosthetic joint infection (PJI). PJI is a serious complication, potentially requiring prolonged antibiotic therapy, multiple surgical debridements, implant removal (explantation), or even amputation in severe cases.
  \item \textbf{Deep Vein Thrombosis (DVT) and Pulmonary Embolism (PE):} Blood clots can form in the deep veins of the leg (DVT) and potentially dislodge, traveling to the lungs (PE), which can be life-threatening. Prophylactic measures (anticoagulants, compression devices) are routinely employed.
  \item \textbf{Wound healing problems:} Delayed wound healing, dehiscence (wound opening), or skin necrosis can occur, particularly in patients with comorbidities like diabetes, obesity, or peripheral vascular disease.
\end{itemize}

\textbf{Procedure-Specific Complications:}
\begin{itemize}
  \item \textbf{Nerve Injury:} Damage to nerves around the knee, most commonly the common peroneal nerve (leading to foot drop or numbness on the top of the foot) or the saphenous nerve (leading to numbness on the medial aspect of the leg), can occur, often transient but sometimes permanent.
  \item \textbf{Vascular Injury:} Although rare, damage to major blood vessels in the popliteal fossa (popliteal artery or vein) can occur during bone resection or implant placement, potentially requiring urgent vascular repair and carrying a risk of limb ischemia.
  \item \textbf{Arthrofibrosis (Joint Stiffness):} Excessive scar tissue formation within the joint can limit knee range of motion, sometimes requiring manipulation under anesthesia (MUA) or further surgical intervention (lysis of adhesions).
  \item \textbf{Implant Loosening or Failure:} Over time, the prosthetic components can loosen from the bone (aseptic loosening) due to wear, osteolysis, or poor fixation, or they may wear out, necessitating revision surgery.
  \item \textbf{Periprosthetic Fracture:} A fracture can occur around the implant during or after surgery, often requiring additional fixation, revision surgery, or prolonged immobilization.
  \item \textbf{Instability:} Ligamentous imbalance, improper implant positioning, or component malrotation can lead to a feeling of instability, recurrent dislocations, or subluxation of the patella or polyethylene insert.
  \item \textbf{Patellar Complications:} If the patella is resurfaced, complications can include patellar fracture, patellar maltracking (leading to pain or instability), patellar clunk syndrome, or avascular necrosis of the patella.
  \item \textbf{Persistent Pain:} While TKA is highly effective for pain relief, a small percentage of patients may experience persistent or unexplained pain despite a technically successful surgery, sometimes due to nerve irritation, complex regional pain syndrome, or other factors.
  \item \textbf{Leg Length Discrepancy:} Although surgeons strive for equal leg lengths, minor discrepancies can occur, which may require shoe lifts or further intervention if symptomatic.
\end{itemize}

\section*{Postoperative Recovery Timeline}
The postoperative recovery timeline for total knee arthroplasty is a structured process designed to maximize functional recovery and minimize complications, typically spanning several months.

Immediately after surgery, the patient is transferred to the Post-Anesthesia Care Unit (PACU) for close monitoring of vital signs, pain levels, and the surgical site. Pain management is a critical priority, often involving a multimodal approach with intravenous medications, oral analgesics, and continued regional nerve blocks. Early mobilization is strongly encouraged, often beginning on the same day or the day after surgery, with assistance from physical therapists. Patients typically start with gentle range-of-motion exercises, partial weight-bearing with a walker or crutches, and exercises to prevent deep vein thrombosis.

The hospital stay usually ranges from 1 to 3 days. During this time, physical therapy sessions intensify, focusing on regaining knee flexion and extension, strengthening the quadriceps and hamstring muscles, and practicing safe transfers and ambulation. DVT prophylaxis (e.g., oral anticoagulants, compression stockings, sequential compression devices) is continued. Upon discharge, patients transition to either home health physical therapy, outpatient physical therapy, or, for some, a short stay in a dedicated rehabilitation facility. The first 1-2 weeks involve managing pain and swelling, gradually increasing activity, and adhering to a prescribed exercise regimen. Most patients can walk with minimal assistance by 4-6 weeks. Full recovery and return to most light activities, such as walking, swimming, and cycling, can take 3 to 6 months, with continued improvement in strength and endurance for up to a year. High-impact activities are generally discouraged long-term.

\section*{Surgical Findings \& Pathology}
During total knee arthroplasty, the surgeon meticulously documents the intraoperative findings, which typically include the extent of articular cartilage loss (often described as eburnated bone or full-thickness defects), the presence and size of osteophytes (bone spurs), the degree of meniscal degeneration or tears, and the severity of any ligamentous laxity or contracture contributing to varus or valgus deformity. The condition of the patellar cartilage and tracking is also assessed. These findings confirm the preoperative diagnosis of end-stage arthritis and guide the precise bone resections and soft tissue releases required for optimal implant placement and knee balance.

The resected bone and cartilage specimens (distal femoral condyles, proximal tibial plateau, and posterior patella) are typically sent for pathological examination. The pathology report usually confirms the presence of degenerative changes consistent with osteoarthritis, such as cartilage fibrillation, erosion, subchondral sclerosis, and osteophyte formation. While the pathological findings are important for confirming the diagnosis, they are generally not used to guide immediate intraoperative decisions, as the diagnosis of severe arthritis is well-established preoperatively through clinical examination and imaging. The primary purpose of the pathology report in TKA is to provide a definitive histological confirmation of the disease process and to rule out any unexpected findings, such as tumor or infection, though these are rare.

\section{Expected Postoperative Course}
A normal and successful postoperative course following total knee arthroplasty is characterized by a progressive and sustained improvement in pain, function, and quality of life. Patients can expect a significant reduction or complete elimination of the severe preoperative knee pain, allowing for improved sleep and reduced reliance on pain medication within weeks to months. Functionally, there should be a steady increase in the ability to walk longer distances, climb stairs, and perform activities of daily living without severe discomfort. The knee joint should feel stable, with a functional range of motion, typically achieving at least 0-110 degrees of flexion, which is crucial for daily tasks. The surgical incision should heal cleanly, without signs of infection, excessive redness, or drainage, usually within 2-3 weeks. Swelling and bruising are common initially but should gradually subside. Ultimately, a normal course means the patient achieves their goals of pain reduction and improved mobility, leading to a better quality of life and a return to most desired low-impact activities within 3-6 months.

\section{Postoperative Care \& Follow-up}
Comprehensive postoperative care and diligent follow-up are essential for ensuring optimal recovery and long-term success after total knee arthroplasty.

\begin{itemize}
  \item \textbf{Postoperative Clinic Visits:} Regular follow-up appointments with the orthopedic surgeon are scheduled, typically at 2 weeks, 6 weeks, 3 months, 6 months, and 1 year post-surgery, and then annually or biennially for long-term monitoring.
  \item \textbf{Suture/Staple Removal:} Skin sutures or staples are usually removed at the first postoperative visit, around 10-14 days after surgery, once the incision has adequately healed.
  \item \textbf{Physical Therapy:} Intensive physical therapy is a cornerstone of recovery, often continuing for 6-12 weeks, and sometimes longer, focusing on strengthening, range of motion, gait training, and functional exercises. Adherence to the home exercise program is crucial.
  \item \textbf{Imaging:} Follow-up X-rays of the knee are typically taken at 6 weeks, 3 months, 1 year, and then periodically (e.g., every 2-5 years) to monitor implant position, alignment, and detect any signs of loosening, wear, or periprosthetic changes.
  \item \textbf{Activity Resumption:} Patients are guided on a gradual return to activities, avoiding high-impact sports but encouraging walking, swimming, cycling, golf, and other low-impact recreational pursuits.
  \item \textbf{Warning Signs:} Patients are educated on critical warning signs to report to their surgeon immediately, such as increasing pain, redness, swelling, warmth, fever, chills, purulent drainage from the incision, or sudden instability or inability to bear weight.
  \item \textbf{Dental Prophylaxis:} Patients are often advised to take prophylactic antibiotics before dental procedures for a period (e.g., 2 years) after surgery to prevent bacteremia from potentially seeding the new joint and causing infection.
\end{itemize}
Long-term care focuses on maintaining joint health, monitoring implant integrity, and addressing any new symptoms.

\section*{Epidemiology}
Total knee arthroplasty is one of the most common and highly successful orthopedic procedures performed globally, with a steadily increasing incidence over the past several decades. This rise is primarily driven by an aging population, the growing prevalence of obesity, and expanded indications for surgery. In the United States, over 700,000 TKAs are performed annually, with projections suggesting a continued substantial increase in demand. The procedure is most common in individuals over the age of 60, although it is increasingly being performed in younger, more active patients. Women tend to undergo TKA more frequently than men, partly due to a higher prevalence of osteoarthritis and longer life expectancy. The primary indication for TKA remains end-stage osteoarthritis, accounting for over 90\% of cases, followed by rheumatoid arthritis and post-traumatic arthritis. Mortality rates associated with elective TKA are remarkably low, typically less than 0.5\%, while major morbidity rates (e.g., deep infection, DVT/PE) range from 1\% to 5\%, underscoring the overall safety and efficacy of this transformative procedure.

\section*{Outcomes \& Prognosis}
The outcomes and prognosis following total knee arthroplasty are overwhelmingly positive, with high rates of patient satisfaction and significant improvements in quality of life. Success rates for pain relief and functional improvement are consistently reported to be over 90-95\%. The vast majority of patients experience substantial reduction or complete elimination of knee pain and a significant increase in their ability to perform daily activities, including walking, climbing stairs, and engaging in recreational pursuits. Implant survival rates are also very favorable, with 10-year survival rates typically exceeding 90-95\% and 20-year survival rates often in the range of 80-85\%, as demonstrated by large national joint replacement registries.

Prognostic factors influencing long-term outcomes include patient age (younger patients may have higher revision rates due to longer implant exposure), activity level, body mass index (obesity is associated with higher complication rates and potentially earlier implant wear), presence of medical comorbidities, and adherence to postoperative rehabilitation protocols. While recurrence of severe arthritis within the replaced joint is not possible, complications such as infection, aseptic loosening, or persistent stiffness can necessitate revision surgery. However, with modern surgical techniques, improved implant designs, and advanced biomaterials, TKA offers a durable and highly effective solution for end-stage knee arthritis, allowing most patients to return to a more active, pain-free, and independent lifestyle for many years.

\section*{References}
\begin{enumerate}
  \item MedlinePlus. Knee Replacement. U.S. National Library of Medicine; 2024. Available from: \url{https://medlineplus.gov/kneereplacement.html}
  \item Insall JN, Scott WN, Ranawat CS. The Total Condylar Knee Prosthesis: A Report of 220 Cases. J Bone Joint Surg Am. 1979;61(2):173-180.
  \item American Academy of Orthopaedic Surgeons (AAOS). Total Knee Replacement. AAOS; 2023. Available from: \url{https://orthoinfo.aaos.org/en/treatment/total-knee-replacement/}
  \item Canale ST, Beaty JH. Campbell's Operative Orthopaedics. 14th ed. Philadelphia, PA: Elsevier; 2021.
  \item Kurtz SM, Ong KL, Lau E, et al. Projections of primary and revision hip and knee arthroplasty in the United States from 2005 to 2030. J Bone Joint Surg Am. 2007;89(4):780-785.
\end{enumerate}

\clearpage